%%%%%%%%%%%%%%%%%%%%%%%%%%%%%%%%%%%%%%%%%%%
\newif\ifinclude		% if per l'include dei file
%\includetrue	

%\newif\ifcomment 	%creo una condizione che mi permette d'usare i commenti multilinea. Andrò a commentare le parti specifiche di ogni file affinche non giri da solo ma assieme a tutto il resto.
%\commentfalse		%Hide comments

%%%%%%%%%%%%%%%%%%%%%%%%%%%%%%%%%%%%%%%%%%%
%
% Template per Elaborato di Laurea
% DISI - Dipartimento di Ingegneria e Scienza dell’Informazione
%
% update 2015-09-10
%
%%%%%%%%%%%%%%%%%%%%%%%%%%%%%%%%%%%%%%%%%%%

% formato FRONTE RETRO
\documentclass[epsfig,a4paper,11pt,titlepage,twoside,openany]{book}
\usepackage{epsfig}
\usepackage{plain}
\usepackage{setspace}
\usepackage[paperheight=29.7cm,paperwidth=21cm,outer=1.5cm,inner=2.5cm,top=2cm,bottom=2cm]{geometry} % per definizione layout
\usepackage{titlesec} % per formato custom dei titoli dei capitoli

\singlespacing

\usepackage[english]{babel}

%%%%%%%%%%%%%%%%% STEFANO ADDED %%%%%%%%%%%%%%%%
% supporto lettere accentate
% queste due righe nel preambolo servono a poter utilizzare le lettere accentate in tutto il testo se no di norma si inserirebbero con \'e...
\usepackage[T1]{fontenc}
\usepackage[utf8]{inputenc}
\usepackage{hyperref} %Serve per i riferimenti
\usepackage{caption}  %Serve per le note (caption)
\usepackage{multicol}	 %Serve per usare più colonne

%%%%%%%%%%%%%%%%% ENRICO ADDED %%%%%%%%%%%%%%%%%
\usepackage{ulem} %Serve per il sottolineato
\usepackage{amsmath} %Serve per alcuni ambienti matematici
\usepackage{array} %Serve per le tabelle
\usepackage{multirow} %Seve per tabelle
\setcounter{secnumdepth}{5} %Utilissimo serve per aumentare il numero di paragrafi, si arriva fino a 5 livelli di profondità x.x.x.x.x
%GRAFICI
\usepackage{pgfplots}
\usepackage{pgfmath}
\usepackage{tikz}
%CODICE
\usepackage{listings}
\usepackage[cache=false]{minted} 
%\setminted{tabsize=4, breaklines, breakanywhere, linenos, mathescape,}

\lstset{
	  breakatwhitespace=false,         
	  breaklines=true,     
	  basicstyle=\footnotesize\ttfamily,            
	  commentstyle=\color{blue}, %Indica il colore dei commenti
	  keywordstyle=\color{red}, %Indica il colore delle parole chiave
	  language=C, %Indica il linguaggio predefinito da usare
	  rulecolor=\color{black}, %Indica il colore dei numeri di righe
	  tabsize=4,
	  escapeinside={\%*}{*)},
	  morekeywords={}, %Altre parole da inserire tra le keywords. Ad esempio possiamo aggiungere do, gotttto, ecc ecc 
}
%%%%%%%%%%%%%%%%% END OF ENRICO  %%%%%%%%%%%%%%%%

\begin{document}
%set the language of the text to english
% !TeX spellcheck = en_GB

%%%%%%% personal commands:
% \newcommand{\nome_commando}[argomenti]{comando}
\newcommand{\mmax}[0]{mod\_withMax}
\newcommand{\mover}[0]{mod\_overlap }
\newcommand{\mmod}[0]{modified modularity}
	
 	% nessuna numerazione
	\ifinclude
  		\pagenumbering{gobble} 
  		\pagestyle{plain}

\thispagestyle{empty}

\begin{center}
  \begin{figure}[h!]
    \centerline{\psfig{file=images/logo_unitn_black_center.eps,width=0.6\textwidth}}
  \end{figure}

  \vspace{2 cm} 

  \LARGE{Dipartimento di Ingegneria e Scienza dell’Informazione\\}

  \vspace{1 cm} 
  \Large{Corso di Laurea in\\
    Informatica\\
    %Informatica
    %Ingegneria dell'Informazione e delle Comunicazioni
    %Ingegneria dell'Informazione e Organizzazione d'Impresa
    %Ingegneria Elettronica e delle Telecomunicazioni    
  }
  \vspace{1 cm} 
  \Large{Percrso di studio\\
    Scienze e Tecnologie Informatiche 
  }
  	

  \vspace{2 cm} 
  \Large\textsc{Elaborato finale\\} 
  \vspace{1 cm} 
  \Huge\textsc{Titolo???\\}
  \Large{\it{Sottotitolo (alcune volte lungo - opzionale)?}}


  \vspace{2 cm} 
  \begin{tabular*}{\textwidth}{ c @{\extracolsep{\fill}} c }
  \Large{Supervisore} & \Large{Laureando}\\
  \Large{Montresor Alberto}& \Large{Leonardi Stefano}\\
  \end{tabular*}

  \vspace{2 cm} 

  \Large{Anno accademico 2017/2018}
  
\end{center}


  		\clearpage
  	 \fi
  
%%%%%%%%%%%%%%%%%%%%%%%%%%%%%%%%%%%%%%%%%%%
%% Nota: Sezione Ringraziamenti opzionale
  	\ifinclude
  		\thispagestyle{empty}

\begin{center}
  {\bf \Huge Ringraziamenti}
\end{center}

\vspace{2cm}
\noindent Ringrazio tutti coloro che in questi anni mi sono stati vicini, mi hanno motivato a crescere e a non mollare anche quando tutto è andato storto. Grazie a chi ha creduto in me, fino a che l'obbiettivo, da lontano e difficile, è diventato sempre più vicino e tangibile.\\
Ora che ho conquistato questo traguardo sono certo che sarà motivo di gioia ed orgoglio per i miei genitori, i parenti e gli amici tutti. Un grazie a tutti gli amici con cui ho condiviso la bellissima esperienza del mondo circense, da cui ho imparato a vivere col sorriso qualsiasi situazione. Altrettanto ringrazio gli amici di più vecchia data, che mi conoscono e m'accompagnano da molti anni.\\
Sono convinto che mi sarete tutti vicini anche per i progetti futuri a cui già sto pensando.\\
Un doveroso ringraziamento va poi a chi ha reso possibile questo lavoro di tesi, il dottorando Nasrullah Sheikh per la sua guida e al prof. Alberto Montresor per la pazienza e i preziosi consigli che mi sono stati di grande aiuto.

  		\clearpage
  		\pagestyle{plain} % nessuna intestazione e pie pagina con numero al centro
  	\fi

  \mainmatter % inizio numerazione pagine in numeri arabi

%%%%%%%%%%%%%%%%%%%%%%%%%%%%%%%%%%%%%%%%%%%
%% Nota:  Si ricorda che il numero massimo di facciate è 30.
%% Nel conteggio delle facciate sono incluse 
%%   indice
%%   sommario
%%   capitoli
%% Dal conteggio delle facciate sono escluse
%%   frontespizio
%%   ringraziamenti
%%   allegati    
%%%%%%%%%%%%%%%%%%%%%%%%%%%%%%%%%%%%%%%%%%%

    % indice
    \ifinclude
    	\tableofcontents
    	\clearpage
    \fi

    % gruppo per definizone di successione capitoli senza interruzione di pagina
    \begingroup
    
    % nessuna interruzione di pagina tra capitoli
    % ridefinizione dei comandi di clear page
    \renewcommand{\cleardoublepage}{} 
    \renewcommand{\clearpage}{} 
    
    % redefinizione del formato del titolo del capitolo
    % da formato
    %   Capitolo X
    %   Titolo capitolo
    % a formato
    %   X   Titolo capitolo
    \titleformat{\chapter}
    	{\normalfont\Huge\bfseries}{\thechapter}{1em}{}
    	
      	\titlespacing*{\chapter}{0pt}{0.59in}{0.02in}
      	\titlespacing*{\section}{0pt}{0.20in}{0.02in}
      	\titlespacing*{\subsection}{0pt}{0.10in}{0.02in}
      
	%\ifinclude
      	% sommario
      	%%%%%%%%%%%%%%%%%%%%%%%%%%%%%%%%%%%%%%%%%%%
\newif\ifcomment	%allow multiline comments
%\commenttrue 		%Show comments
%%%%%%%%%%%%%%%%%%%%%%%%%%%%%%%%%%%%%%%%%%%
\ifcomment	%<<<<<<<<<<<<<<<<<<<<<<<<<<<<<<<<<<<<<<<<<<<<<<<<<<<<<<<<<
	\documentclass[epsfig,a4paper,11pt,titlepage,twoside,openany]{book}

	% queste due righe nel preambolo servono a poter utilizzare le lettere accentate in tutto i testo se no 	di norma si inserirebbero con \'e...
	\usepackage[T1]{fontenc}
	\usepackage[utf8]{inputenc}
	
	\usepackage{hyperref} %Serve per i riferimenti
	\usepackage{caption}  %Serve per le note (caption)
	
	\begin{document}
	%set the language of the text to english
	% !TeX spellcheck = en_GB
	
	% personal commands:
	% \newcommand{\nome_commando}[argomenti]{comando}
	\newcommand{\mmax}[0]{mod\_withMax}
	\newcommand{\mover}[0]{mod\_overlap }
	\newcommand{\mmod}[0]{modified modularity}
\fi					%<<<<<<<<<<<<<<<<<<<<<<<<<<<<<<<<<<<<<<<<<<<<<<<<<<<<<<<<<
%
\chapter{Summary}
% Montresor suggerisce di scriverlo alla fine in quanto solo allora si avrà una visione d'insieme di tutto il lavoro.
%% cos'è?
%% Sommario è un breve riassunto del lavoro svolto dove si descrive: l’obiettivo, l’oggetto della tesi, le metodologie e le tecniche usate, i dati elaborati e la spiegazione delle conclusioni alle quali siete arrivati.
%% Il sommario dell’elaborato consiste al massimo di 3 pagine e deve contenere le seguenti informazioni: contesto e motivazioni breve riassunto del problema affrontato tecniche utilizzate e/o sviluppate risultati raggiunti, sottolineando il contributo personale del laureando/a se vi sono più persone che hanno contribuito...
This thesis goes back over the work that I was done during the internship inside the university, with the supervision of the PhD student Sheikh Nasrullah for the professor Montresor Alberto.
%
\section{Introduction}
The whole internship is based on the paper named "Community-enhanced Network Representation Learning for Network Analysis" alias CNRL, this paper proposes an innovative method for detecting the communities inside a graph, so my purpose is to take the code at the base of this paper, try to replicate the results shown and after that increase if it is possible the potentiality of this algorithm.\newline
At the start, I had to need to understand the main problem and understand the dynamics at the base of it, therefore I read out some other papers like:
\begin{itemize}
	\item "DeepWalk: Online Learning of Social Representations"
	\item "node2vec: Scalable Feature Learning for Networks"
	\item "How exactly does word2vec work?"
\end{itemize} 
Before the explaining of the work done, we need to understand why the community detection is so important.
%
\section{What are graphs and the communities}
At the base of all, there are the graphs, a particular data structure formed of the nodes that can represent an entity, eventually with some attributes, and the edges that link two nodes, the edges are all directed or not. With this special data structure, we can represent a huge amount of situations. The purpose of the community detection is to find the communities inside the graph, that are groups of nodes that are similar.\\
What does this mean?\\
Two nodes are similar if have an equivalent role inside a group of nodes, they can are a hub of connection or an isolated point or more simply they have the same value for some attribute. To find this groups can permit to who look for it, to manage those nodes in a similar way. Like the possible application of the graphs, the same is for this technology, there are lots of situation and context where those algorithms can work very well.
%
\section{How we evaluate the community}
Up to now, I have explained what are graphs and communities, and from my words seem that when you choose some groups of nodes is may be easily understood if it is a good partition or not. And if you have two partitions to decide which is the best is not difficult. This is not really true. If we have a little graph with only one community we can look at it and decide if the partition is good or not, but this is not possible for a big graph.\\
So how we evaluate a partition?\\
First of all, we must explain what we intend with partitions, here the characteristics:
\begin{itemize}
	\item A partition is a set of communities
	\item Not all the communities have the same size, or better the same number of nodes
	\item For the communities, there is not a minimum size neither a maximum size
	\item Is possible that some nodes do not have a community to which it belongs
	\item A node can belong to more than one community, this mean that I have an overlapping of communities
\end{itemize}
Said that, we can speak about the modularity.\\
The modularity is the method used for evaluating a partition, there are different implementations of it because the aim is the same but we can have different interest so we can weight different aspect, in our work we have used three different methods: \mmax, \mover and \mmod . The first two are unofficial names, we have chosen them for distinguishing it, the last one is the method used in the CNRL's paper. More explanation of those methods in the chapter on the modularity.
%
\section{Our changes}
The code of CNRL do $x$ random walks (a sequence of $l$ adjacent nodes) on the graph where $x=n*w$: 
\begin{itemize}
	\item $n$: is the number of nodes in the graph
	\item $w$: is the number of walks starting from each node
	\item $l$: is the maximum length of each walk
\end{itemize}
Given those walks, the algorithm calculates the partition. Note that CNRL use only the structure of the graph, the nodes and the edges, no other information is considered.\\
Our approach uses the code of CNRL but introduces the attributes of the nodes, this permit a better evaluation. This because two nodes that before are far, there was no one link between them, now are considered near if they share one or more attributes.
%
\section{Results}
To be continued...

%
%
%
%
\newpage
%
\ifcomment	%<<<<<<<<<<<<<<<<<<<<<<<<<<<<<<<<<<<<<<<<<<<<<<<<<<<<<<<<<
	\end{document}
\fi					%<<<<<<<<<<<<<<<<<<<<<<<<<<<<<<<<<<<<<<<<<<<<<<<<<<<<<<<<<

 	%\fi
      
%%%%%%%%%%%%%%%%%%%%%%%%%%%%%%%%%%%%%%%%%%%
%% Nota:
%% Sommario e' un breve riassunto del lavoro svolto dove si descrive 
%% l’obiettivo, l’oggetto della tesi, le metodologie e 
%% le tecniche usate, i dati elaborati e la spiegazione delle conclusioni 
%% alle quali siete arrivati.
%% Il sommario dell’elaborato consiste al massimo di 3 pagine e deve contenere le seguenti informazioni: 
%%   contesto e motivazioni
%%   breve riassunto del problema affrontato
%%   tecniche utilizzate e/o sviluppate
%%   risultati raggiunti
%%%%%%%%%%%%%%%%%%%%%%%%%%%%%%%%%%%%%%%%%%% 
      
	%%%%%%%%%%%%%%%%%%%%%%%%%%%%%%%%
    % lista dei capitoli
    %
    % \input oppure \include
    %
    	\input{1_Implementation}		%capitolo 1
      	%%%%%%%%%%%%%%%%%%%%%%%%%%%%%%%%%%%%%%%%%%%%%%
\newif\ifcomment	%allow multiline comments
%\commenttrue 		%Show comments
%%%%%%%%%%%%%%%%%%%%%%%%%%%%%%%%%%%%%%%%%%%%%%
\ifcomment	%<<<<<<<<<<<<<<<<<<<<<<<<<<<<<<<<<<<<<<<<<<<<<<<<<<<<<<<<<
\documentclass[epsfig,a4paper,11pt,titlepage,twoside,openany]{book}

% queste due righe nel preambolo servono a poter utilizzare le lettere accentate in tutto i testo se no di norma si inserirebbero con \'e...
\usepackage[T1]{fontenc}
\usepackage[utf8]{inputenc}

\usepackage{amsmath}
\usepackage{caption}  %Serve per le note (caption)
\usepackage{hyperref} %Serve per i riferimenti

\begin{document}
%set the language of the text to english
% !TeX spellcheck = en_GB

% personal commands:
% \newcommand{\nome_commando}[argomenti]{comando}
\newcommand{\mmax}[0]{mod\_withMax}
\newcommand{\mover}[0]{mod\_overlap }
\newcommand{\mmod}[0]{modified modularity}
\fi					%<<<<<<<<<<<<<<<<<<<<<<<<<<<<<<<<<<<<<<<<<<<<<<<<<<<<<<<<<

%
\chapter{Evaluation methods}
In this section, I am going to explain how we evaluate the partitions created by the community detection algorithms, as I introduce in the summary, this method, that we decided to use, is called modularity. Its symbol is $Q$.\\
The modularity is a value that is typically contained in the range of $[-1, 1]$, but this is not all times true because it depends on the implementations. The first persons who have worked on the modularity were Newman and Girvan, that like us, they had needed a method for choosing between two partitions of communities on the nodes.\\
They start from an assumption a node can not belong to more than one community, this assumption is not always true for our situation but it simplifies a lot the complex problem of calculating modularity. So for the first steps, we use the same assumption.\\
During the time of the internship we use three different methods, not ideate from our, now I am going to explain them from the simplest to the most complex. The formulas reported are for undirected graphs at the end of each subsection there is the differences with directed graphs.
%
%
\section{Modularity with Maximum}
We named it in this method because it permits at maximum one community for each node, perfectly in line with the assumption of Newman and Girvan.\\
The idea is that this method calculates a modularity value for each community in the graph and the total modularity is simply the summation of all the parts. We must specify that, in this method, the overlapping is not considered then the summation does not consider the same elements more than one times.\\
This method is based on the observation that we can identify a community on the base of the frequency of its edges, in fact, if a set of points have a huge amount of edges inside it probably this is a community if it is not, this is a bad situation that must penalize all the partition.\\
We chose it because is easy to calculate and because it penalizes a lot the wrong partitions.\\
Here the formula explained:
\begin{equation}
	Q=\sum_{c}^{C_r} \left( \frac{l_c}{L}-\left(\frac{d_c}{2L} \right)^2 \right)
	\label{eq:m_max}
\end{equation}
Where the elements of the formula are:
\begin{itemize}
	\item $Q$ is the symbol for modularity
	\item $C_r$ is the set of the communities
	\item $c$ is the iterator on the communities
	\item $L$ is the total number of edges in the graph
	\item $l_c$ is the total number of edges inside, both nodes of each edge are inside the community $c$
	\item $d_c$ is the degree of the community $c$, defined as the summation of the degree of the nodes belong to it
\end{itemize}
We can see that $\displaystyle\frac{l_c}{L}$ is the weight of the community compared to the others in the graph, and $\displaystyle\frac{d_c}{2L}$ is the density of edges, in other words, the percentage of edges that belong to this community compared to the total number.\\
\\
To note that when the graph is directed change only a multiplicative factor indeed the "density of edges" change from $\displaystyle\frac{d_c}{2L}$ to $\displaystyle\frac{d_c}{L}$, and this is all.
%
%
\section{Modularity with Overlap}
This algorithm is taken from the paper named "Modularity measure of networks with overlapping communities".\\
As the name says, now the overlapping of the community it is allowed. Then now we do not have a summation but an average. The results are inside a range of $[-1, 1]$.\\
Here the formula explained:
\begin{equation}
	Q = M^{ov} = \frac{1}{K} 
	\sum_{r=1}^{K} \left[
		\frac
			{\sum\limits_{i \in c_r} 
				\left( \frac
					{
						\sum\limits_{j \in c_r, i \neq j} \left( a_{ij} \right) 
						- 
						\sum\limits_{j \notin c_r} \left( a_{ij} \right) 
					} 
					{d_i \cdot s_i} 
				\right) } 
			{n_{c_r}}
		\cdot
		\frac{ n^e_{c_r} }{ \binom{n_{c_r}}{2} } 
	\right]
	\label{eq:m_over}
\end{equation}
Where the elements of the formula are:
\begin{itemize}
	\item $K$ is the number of communities $|C_r|$
	\item $c_r$ is the actual community
	\item $n_{cr}$ number of nodes in the community
	\item $n^e_{c_r}$ number of edges in the community
	\item $\binom{n_{c_r}}{2}$ maximum number of edges in the community
	\item $d_i$ degree of the node $i$
	\item $s_i$ number of community to which the node $i$ belong
	\item $a_{ij}$ 1 if the edge from the node $i$ to the node $j$ exist, 0 otherwise
\end{itemize}
Now, we try to figure out the particular things that appear in this formula.\\
As before we calculate the density of the community (with the equation number~\ref{eq:edge_density}), but not compared to the total number of edges of the graph, but compared to the hypothetical maximum number of edges inside the nodes considered:
\begin{equation}
	\frac{ n^e_{c_r} }{ \binom{n_{c_r}}{2} }
	\label{eq:edge_density}
\end{equation}
In addition, we consider the relationship between the number of inward and outward edges, compared to the total number of edges of the community (in the equation~\ref{eq:in-out_edge_relation}). Then if the outward edges are more than the inward edges, the modularity is negative otherwise is positive.\\
\begin{equation}
	\sum\limits_{i \in c_r} \left( \frac{
		\sum\limits_{j \in c_r, i \neq j} \left( a_{ij} \right) - 
		\sum\limits_{j \notin c_r} \left( a_{ij} \right) 
	} {d_i \cdot s_i} \right)
	\label{eq:in-out_edge_relation}
\end{equation}
Lastly, when we look at the inward and outward edges we consider the two parameters $d_i$ and $s_i$, this is important because if a node has only one edge and it belongs only to one community it has a high weight, differently if it has a hundred of nodes and belongs to dozens of communities, the weight of one of its edges is very low.\\
\\
We chose this method because it allows the overlapping of communities and this is necessary for our situation. In addition is it possible to figure out which community have a high value, then increase the final results and which one decrease it, cause the low value.\\
From each modularity value, we can understand two things of the graph:
\begin{itemize}
	\item If the value is negative mean that the group of nodes is not densely connected inside it, but rather is connected to some external point. It is possible that there are more inwards edges compared to the outwards, but if the value is negative this means that the outwards edges are more relevant
	\item If the modularity in absolute value is very near at 0 this means that the graph is extremely sparse in fact the multiplication factor (shown in equation~\ref{eq:edge_density}) collapse all the elements of the formula near zero
\end{itemize}
\noindent To note that when the graph are directed only the equation~\ref{eq:edge_density} change in fact the maximum number of edges is no more $ \displaystyle\binom{n_{c_r}}{2}$ but became $\displaystyle 2\binom{n_{c_r}}{2}$.
%
%
\section{Modified modularity}
This algorithm is taken from the paper named "Incorporating Implicit Link Preference Into Overlapping Community Detection".\\
Finally, I arrive to explain this method, this is the algorithm choose in the paper of CNRL, for this reason, we looked for it. It is very important because allow our to replicate the results shown in the paper of CNRL, then we can proceed to find a fixed point. Unfortunately, this is the slowest algorithm of the three that I show you in this section because it is the most complex.\\
Let's look at how it works:
\begin{equation}
	Q = \frac{1}{M} \cdot \sum_{u,v \in V}
		\left(
			\left( A \left[ u,v \right] - \frac{ d_{in}\left(u\right) \cdot d_{out}\left(v\right) }{M} \right)
			\cdot
			|C_u \cap C_v| 
		\right)
	\label{eq:m_mod}
\end{equation}
Where the elements of the formula are:
\begin{itemize}
	\item $M$ stands for $m$ if the graph is undirected and for $2m$ if the graph is directed, where $m$ is the number of edges of the graph
	\item $u,v$ are the iterators on the graph's nodes named $V$
	\item $A \left[ u,v \right]$ is the weight of the edge $uv$ if the edge does not exist this value is 0, normally a node without weight is considered as 1. It is taken from the adjacency matrix
	\item $d_{in}\left(u\right)$ and $d_{out}\left(u\right)$ are the incoming/outgoing degree namely the number of incoming/outgoing edges of the node $u$. If the graph is undirected those values are equal
	\item $C_u$ is the set of communities to which the node u belong
	\item $|C_u \cap C_v|$ is the number of communities that node $u$ and $v$ share
\end{itemize}
I try to explain this equation in words.\footnote{Non sono completamente sicuro delle meccaniche alla base della modularità modificata}\\
Starting from the end we see the element $\displaystyle |C_u \cap C_v|$, this is important because allow ignoring the edges that link two nodes of different communities, so only the edges inside a community is considered. In addition, this part is a multiplicative factor, if an edge is inside lots of communities it has a high weight if it belongs only to one community is weight is lower.\\
But what is the base of this multiplier factor? We can see that this is a positive value if the edge exists (thanks to $\displaystyle A \left[ u,v \right]$), in fact, more edges a community has more it is defined. From the weight of the edge we subtract a factor (defined from $\displaystyle \frac{ d_{in}\left(u\right) \cdot d_{out}\left(v\right) }{M}$ ), this represents the relevance/likelihood of the edge, if it connects two nodes with a high degree, this edge is not too important therefore we decrease its weight slightly, at the opposite is very important if connect two nodes that have low values of degree.
%
%
%
%
%
%
\ifcomment	%<<<<<<<<<<<<<<<<<<<<<<<<<<<<<<<<<<<<<<<<<<<<<<<<<<<<<<<<<
\end{document}
\fi					%<<<<<<<<<<<<<<<<<<<<<<<<<<<<<<<<<<<<<<<<<<<<<<<<<<<<<<<<<
	%capitolo 2
      	%\input{capitolo3}	%capitolo 3
      	%\input{capitolo4}	%capitolo 4
    \endgroup


    % bibliografia in formato bibtex
    % aggiunta del capitolo nell'indice
    \addcontentsline{toc}{chapter}{Bibliografia}
    % stile con ordinamento alfabetico in funzione degli autori
    \bibliographystyle{plain}
    \ifinclude
    	%\bibliography{biblio}
    \fi
%%%%%%%%%%%%%%%%%%%%%%%%%%%%%%%%%%%%%%%%%%%
%% Nota:
%%
%% Nella bibliografia devono essere riportati tutte le fonti consultate 
%% per lo svolgimento della tesi. La bibliografia deve essere redatta 
%% in ordine alfabetico sul cognome del primo autore. 
%% 
%% La forma della citazione bibliografica va inserita secondo la fonte utilizzata:
%% 
%% LIBRI
%% Cognome e iniziale del nome autore/autori, la data di edizione, titolo, casa editrice, eventuale numero dell’edizione. 
%% 
%% ARTICOLI DI RIVISTA
%% Cognome e iniziale del nome autore/autori, titolo articolo, titolo rivista, volume, numero, numero di pagine.
%% 
%% ARTICOLI DI CONFERENZA
%% Cognome e iniziale del nome autore/autori (anno), titolo articolo, titolo conferenza, luogo della conferenza (città e paese), date della conferenza, numero di pagine. 
%% 
%% SITOGRAFIA
%% La sitografia contiene un elenco di indirizzi Web consultati e disposti in ordine alfabetico. 
%% E’ necessario:
%%   Copiare la URL (l’indirizzo web) specifica della pagina consultata
%%   Se disponibile, indicare il cognome e nome dell’autore, il titolo ed eventuale sottotitolo del testo
%%   Se disponibile, inserire la data di ultima consultazione della risorsa (gg/mm/aaaa).    
%%%%%%%%%%%%%%%%%%%%%%%%%%%%%%%%%%%%%%%%%%%
%%%%%%%%%%%%%%%%%%%%%%%%%%%%%%%%%%%%%%%%%%%    

    \titleformat{\chapter}
        {\normalfont\Huge\bfseries}{Allegato \thechapter}{1em}{}
    % sezione Allegati - opzionale
    \appendix
    %\input{allegati}

\end{document}
